\clearpage
\VerdierPhysicalPageAnchor{0218}%
\begingroup\small\setlength{\parskip}{0pt}\setlength{\emergencystretch}{1.5em}
\setlength{\abovedisplayskip}{4pt plus 2pt minus 2pt}%
\setlength{\belowdisplayskip}{4pt plus 2pt minus 2pt}%

\VerdierSemanticAnchor{ch3d-3-3-1-2}
\noindent\textit{$S_i^1*f_{i-1}$、$1<i\leq n$、
$f_i*{S''}_i^1$、$1\leq i<n$、${S'}_i^1*g_{i-1}$、$1<i\leq n$、
$g_i*{S''}_i^1$、$1\leq i<n$ であり、これらが ((3.2.8.1) と
(3.2.8.2)) により定義され、次を満たすものとする：}

\begin{equation}
\left\{
\begin{aligned}
\delta^1({S''}_n^1)
  &=\delta^1(S_n^1*f_{n-1}),\\
\delta^1({S''}_n^1)
  &=\delta^1({S'}_n^1*g_{n-1}),\\
\delta^1(f_i*{S''}_i^1)
  &=\delta^1(S_i^1*f_{i-1}),&&1<i\leq n-1,\\
\delta^1(g_i*{S''}_i^1)
  &=\delta^1({S'}_i^1*g_{i-1}),&&1<i\leq n-1,\\
\delta^1(f_1*{S''}_1^1)
  &=\delta^1(S_1^1),\\
\delta^1(g_1*{S''}_1^1)
  &=\delta^1({S'}_1^1).
\end{aligned}
\right.
\tag{3.3.1.2}
\end{equation}

\VerdierSemanticAnchor{ch3d-3-3-1-proof}
実際、主張 a) は (3.2.2) と (3.2.5) から直ちに従い、主張 b) は写像
$E_{\mathcal A}^{n}(X,Y)\longrightarrow
\operatorname{Ext}_{\mathcal A}^{n}(X,Y)$ が誘導する同値関係の記述
（(3.2.2) と (3.2.9)）から従う。

\VerdierSemanticAnchor{ch3d-3-3-2}
\medskip\noindent\textbf{3.3.2.}\quad
(1.2.11) を思い出すと、圏 $\mathcal A^{\operatorname{Ext}}$ は
$D(\mathcal A)$ の充満部分圏であり、$\mathcal A$ の対象をシフトした
対象によって定まる。
$F\colon\mathcal A^{\operatorname{Ext}}\longrightarrow\mathcal C$ を関手
（値をとる圏は加法圏とは限らない）とする。各整数 $n$ に対し、
$F^n\colon\mathcal A\longrightarrow\mathcal C$ を関手
$X\longmapsto F(X[n])$ とする。同様に、$\mathcal A$ の任意の完全列
$S^1=0\longrightarrow X\longrightarrow Y\longrightarrow Z
\longrightarrow0$ に対して、
$\delta_F^n(S^1)\colon F^n(Z)\longrightarrow F^{n+1}(X)$ を次の射とする：

\VerdierSemanticAnchor{ch3d-3-3-2-delta-definition}
\[
F\bigl(\delta^1(S^1)[n]\bigr)
\colon F(Z[n])\longrightarrow F(X[n+1]).
\]

\VerdierSemanticAnchor{ch3d-3-3-2-sfl}
関手列 $F^n$、$n\in\mathbb Z$ は射 $\delta_F^\ast$ とともに、次の性質を
もつ (3.2.9)：

\noindent\textbf{SFL I)} 任意の可換図式：

\VerdierSemanticAnchor{ch3d-3-3-2-sfl-i-source-diagram}
\begin{center}
\resizebox{0.86\textwidth}{!}{%
$\begin{tikzcd}[ampersand replacement=\&,column sep=3.0em,row sep=2.8em]
S=0 \arrow[r]
  \& X \arrow[r] \arrow[d,"f"']
  \& Y \arrow[r] \arrow[d]
  \& Z \arrow[r] \arrow[d,"g"]
  \& 0 \\
{S'}=0 \arrow[r]
  \& X' \arrow[r]
  \& Y' \arrow[r]
  \& Z' \arrow[r]
  \& 0
\end{tikzcd}$%
}
\end{center}
\endgroup
